% Subsection: Union, Intersection, Complement, and Difference  (notes stub; content authored later)
\subsection{Union, Intersection, Complement, and Difference}

\begin{remark*}[Exposition]
The primitive set operations act first on individual subsets of a fixed
ambient set \(X\).  In this chapter the inputs are usually not arbitrary
subsets, but members of a selected family \(\mathcal{F}\subseteq\mathcal{P}(X)\).
The output is still a subset of \(X\), even before one knows whether the output
belongs to \(\mathcal{F}\).
\end{remark*}

\begin{definition}[Operation on Members of a Family]\label{def:operation-on-members-of-family}
Let \(X\) be a set and let \(\mathcal{F}\subseteq\mathcal{P}(X)\).  An
\emph{operation on members of \(\mathcal{F}\)} is an ordinary set-theoretic
operation whose inputs are chosen from \(\mathcal{F}\) and whose output is
regarded as a subset of \(X\).
\end{definition}

\begin{remark*}[Standard quantified statement]
For a binary operation \(\star\) on subsets of \(X\), applying \(\star\) to
members of \(\mathcal{F}\) means
\[
  A,B\in\mathcal{F}
  \quad\Longrightarrow\quad
  A\star B\subseteq X.
\]
This statement records only that the operation is meaningful inside the
ambient power set.  It does not assert \(A\star B\in\mathcal{F}\).
\end{remark*}

\begin{remark*}[Interpretation]
The distinction between ``the operation is defined'' and ``the family is
closed under the operation'' is essential.  If \(A,B\in\mathcal{F}\), then
\(A\cup B\), \(A\cap B\), and \(A\setminus B\) are always subsets of \(X\).
The next question is whether those resulting subsets remain inside the
selected family.  That is a closure question, not a definition-of-operation
question.
\end{remark*}

\begin{remark*}[Examples]
The operations most frequently applied to members of a family are union,
intersection, complement relative to \(X\), set difference, and symmetric
difference:
\[
  A\cup B,\qquad
  A\cap B,\qquad
  A^c=X\setminus A,\qquad
  A\setminus B,\qquad
  A\triangle B.
\]
\end{remark*}

\begin{dependencies}
\begin{itemize}
  \item \hyperref[def:family-of-subsets]{Family of Subsets}
  \item \hyperref[def:union]{Union}
  \item \hyperref[def:intersection]{Intersection}
  \item \hyperref[def:complement]{Complement}
  \item \hyperref[def:set-difference]{Set Difference}
\end{itemize}
\end{dependencies}
